\chapter{INTRODUCTION}

Music theory is the analysis of music through different aspects, such as melody and harmony. Music theory has shown to be useful in both music composition and performance as it allows for a more structured analysis of the music piece \cite{OrioRoda}. Furthermore, music theory is deeply interconnected with musicology as musicologists must be aware of how to analyze music in order to properly understand and give context for it \cite{Burkholder}.

Representing music through a graph-based representation is a novel method of performing this analysis. Aside from capturing distinct segments of a musical piece, it allows the relationships between these segments throughout the piece to be preserved and observed \cite{HernandezOlivanBeltran}. Graph-based representations allow for a variety of analyses to be performed, such as analyses of the structure, melody, and harmony of the piece \cite{ClarkCook, ConananEtAl}.

Some approaches to mapping musical pieces into data, such as the graph-theoretic approach introduced by Conanan, Santuyo and Bomediano, result in sequences of data, or time series, that differ in length \cite{ConananEtAl}. As such, measuring the similarity between these sequences requires an algorithm that can account for one-to-many and many-to-one relationships. Dynamic Time Warping (DTW) is one such algorithm. Using dynamic programming, DTW is able to measure the similarity between two sequences of data points, regardless of whether they are of the same size or not \cite{ConananEtAl}. This is particularly needed in creating distance matrices for graph-based representations of musical compositions.


\section{Context of Study}

Graph representations of musical pieces allow for a novel analysis of the piece with a range of applications \cite{OrioRoda, ClarkCook, HernandezOlivanBeltran}. In 2009, Orio and Roda describe a measure of comparing melodic similarity in music through a graph-based representation. The data used was in the form of MIDI files. The melodic content is simplified and converted into a simpler profile that should retain the most important information, resulting in a single note. The first step in  this graph-based representation is to perform harmonic analysis; this was done manually in this paper. The second step is to simplify the piece; this is done by assigning three weight coefficients to each note: harmonic weight, metric weight, and melodic weight. More relevant notes will have higher weights. The simplification is done over a sliding window until the window reaches the end of the melody. This is done iteratively until the remaining result is one or two notes. These segments are then inserted into a graph where the terminal nodes are the surface-level of the melody and the internal nodes are the generalized versions. The similarity is based on the shortest path between terminal nodes being compared \cite{OrioRoda}.

In 2018, Orio and Roda with Simonetta and Carnovalini improved upon this measure. The dataset was changed from MIDI files to MusicXML files due to the higher amount of information stored in MusicXML files. The paper intends to represent both a contrapuntal and harmonic point of view compared to the previous iteration focusing only on the melodic. The steps from the previous paper of segmentation and reduction are copied from the previous iteration. The reduction step was improved to account for tuplets, triplets, and tied notes in order to improve the accuracy of the representation. Weights were also added to the reduction step in order to account for the difference of the original notes from the reduced notes. Afterwards, an additional step to identify the transformations between the segments is done. In order to measure similarity, the weights of each edge between the terminal nodes being compared are added as opposed to the unweighted edges in the previous paper \cite{SimonettaEtAl}.

Alternatively, Hernandez-Olivan and Beltran proposed two algorithms intended for music boundary detection wherein a piece would be segmented based on its structure as compared to the similarity comparison of the previous papers. These two algorithms, G-PELT and G-Window, were graph-based. Therefore, a graph representation of the piece was created before applying the algorithms. In order to convert the piece into music, they took into account the notes and the time signatures. In this graph representation, they represented three types of edges between the nodes: notes that occur on the same onset, consecutive notes in time, and overlapping notes in time. Afterwards, G-PELT and G-window were used on the graph representation in order to segment the piece. Notably, the G-PELT and G-window algorithms have time complexities of $O(n^2) + O(n^2+nlog(n))$ and $O(n^2)+O(2nw)$ respectively where n is the number of notes and w is the window size \cite{HernandezOlivanBeltran}.

In 2025, Bomediano, Conanan, and Santuyo created a model that takes MusicXML files of classical music compositions and converts it into a node graph that is analyzed for its structure and similarity, following the Implication-Realization Temporal-Gestalt framework. The Implication-Realization (I-R) model developed by Narmour states that listeners have expectations for the melody depending on previous musical events \cite{FosterNarmour}. Tenney and Polansky define temporal-gestalt (TG) units wherein events are hierarchically ordered based on proximity in time and similarity. These TG units can form events, clangs, and sequences wherein similar TG units occur closely or segregation wherein TG units are dissimilar or occur far from each other in time \cite{TenneyPolansky}.

Each composition was converted into a note matrix. The note matrix is then annotated with an I-R pattern, following the rules described by Wen et al \cite{WenEtAl}. The matrix is then segmented wherein notes are grouped based on pitch and duration, following the temporal-gestalt framework. Afterwards, graph construction is performed on the segments \cite{ConananEtAl}.

In the graph, nodes are created from segments of the matrix representation of the musical composition, and edges are created based on the distances between the segments. The entries of the adjacency matrix of the graph is calculated by applying Dynamic Time Warping (DTW) on two segments in order to compute the optimal distance between the note matrices. The nodes and segments are then graphed using the K-Nearest Neighbors model. Lastly, each node is given a label based on the expectancy score and the I-R symbol. After graph construction, the graph identity is evaluated in order to measure its performance, similarity score, and other metrics \cite{ConananEtAl}.

Dynamic Time Warping was used to construct the adjacency matrix because it is able to compare time series of different lengths, a necessity given that the musical segments vary in length. Dynamic Time Warping is also able to capture shape similarity in data across temporal variations \cite{paparrizos}, which is necessary due to the variations that arise from the segmentation of the piece.   

Notably, the standard implementations of DTW have a runtime of $O(n^2)$, therefore the runtime of the graph construction by Conanan et al is bottle-necked by a time complexity of $O((mn)^2)$ due to calculating the distance between every pair. In order to save  time, only the top diagonal of the distance matrix is calculated since the distance matrix is symmetric across the diagonal \cite{ConananEtAl}. Improvements to its time complexity have been researched such as by Gold and Sharir wherein an implementation of DTW with less than $O(n^2)$ time complexity was created \cite{OmerGold}.

Parallelization has been another method of improving the runtime of DTW. Parallelization techniques aim to either reduce the quadratic memory needed or the quadratic computation time taken of DTW. Different studies parallelize DTW in different ways. However, these different methods can be grouped into two categories: algorithm-based parallelization, which restructures the algorithm in order to reduce the memory usage, such as the divide-and-conquer algorithm proposed by Tralie and Dempsey \cite{TralieDempsey}; and hardware-based parallelization, which utilizes the capabilities of different pieces of hardware to speed up the computation of DTW, such is shown by Campos et al \cite{GPUParallelization} wherein different hardware ran the same implementation of DTW and resulted in differing speeds, all of which were faster than the base DTW .






\section{Research Questions}

This study investigates the potential methods of improving the runtime of DTW through implementations of DTW using deterministic algorithm and parallelized implementations of DTW and the extent of these improvements. Specifically, the research addresses the following questions:
\begin{enumerate}
    \item How can deterministic algorithms affect the performance of DTW in graph construction?
    \item To what extent does parallelization affect the runtime of DTW in graph construction?
\end{enumerate}
\section{Research Objectives}

This paper aims to improve the runtime of DTW and apply it to the model for IR-TG graphs created by Conanan et al. These changes are aimed to reduce the runtime of the DTW algorithm and implement parallelization. It seeks to evaluate the extent of these improvements and put them into a single model. Lastly, it aims to show that the improved model has a comparable performance to the original in terms of accuracy. The specific objectives are:
\begin{enumerate}
    \item To implement a version of DTW with a runtime of less than $O(n^2)$ as described by Gold and Sharir.
    \item To implement a parallelized version of the modified DTW so that more than a single calculation of DTW may be done at a single time.
    \item To combine these two improvements into a single model that has a lower runtime than the model by Conanan et al and produces comparable scores in the performance metrics.
\end{enumerate}


\section{Scope and Limitations of the Study}

The framework of node graphs will follow the Implication-Realization Temporal-Gestalt (IR-TG) Graph framework used by Conanan et al in their model \cite{FosterNarmour, TenneyPolansky}. For the consistency and comparison of results, the same dataset of 39 compositions from Bach, Chopin, and Ysaye will be used in this paper as used by Conanan et al \cite{ConananEtAl}. The compositions will be represented as MusicXML and MIDI files and analyzed through the music21 Python library. All of the compositions are monophonic.

The data gathering, feature extraction, segmentation, and performance metrics for graph analysis will be retained from Conanan et al as defined in their methodology. This paper will focus on improving the graph construction portion of the model due to the high runtime creating a bottleneck. The implementation of dynamic time warping used in the model will be modified to potentially have a lesser runtime than the implementation used by Conanan et al. Performance metrics for runtime of the graph will be added to the model alongside the existing performance metrics for graph analysis.
\section{Significance of the Study}

Graph-based representations allow for a novel analysis of music pieces as it allows for the relationships between segments throughout the piece to be captured. The structure of music pieces may go beyond the individual music events as a holistic, nonlinear structure is formed among the events \cite{HernandezOlivanBeltran}. 

Reducing the runtime of this model will allow for increased speed of analysis of music pieces as it will require less resources. This will allow for more comparisons, allowing for improved analysis of pieces as there can be a larger number of comparisons that can be made and increasing the value of each individual comparison \cite{ClarkCook}. Furthermore, this reduces the need of the model on modules in creating checkpoints, such as pickle, thus reducing the potential risk of data loss and wasted computation time. A faster implementation of the model also increases its viability for industry-level applications. 

In addition to the improved similarity comparisons through the graph-representations, DTW is a popular algorithm for many purposes. Exploring a faster implementation of this algorithm and its limitations can improve applications of DTW in many other models, not limited to computational musicology. Furthermore, it allows for further research on DTW and its applications.

